\par\medskip\noindent
\textbf{\textit{Notations and reminders 2.5.0.}}---%
\footnote{Editor's note: this paragraph of notations and reminders has
been added.}
Let \(S\) be a scheme, let \(G\) be an \(S\)-group scheme, and let
\(\epsilon\colon S\to G\) be the unit section of \(G\).  We saw in
II, §\SGARefLink{sga3:xix:paragraph:4-11}{4.11} that the functor
\(\underline{\operatorname{Lie}}(G/S)\) is represented by the vector
bundle
\(\mathbb V(\omega^1_{G/S})\), where
\(\omega^1_{G/S}=\epsilon^{*}(\Omega^1_{G/S})\), and we write
\[
  \mathfrak g
  =
  \operatorname{Lie}(G/S)
  =
  \bigl(\omega^1_{G/S}\bigr)^{\vee}
\]
for the sheaf of sections of this vector bundle.  Suppose in addition
that \(G\) is smooth over \(S\).  Then \(\omega^1_{G/S}\), and hence
\(\operatorname{Lie}(G/S)\), are finite locally free
\(\mathcal O_S\)-modules, and (cf. I, \SGARefLink{sga3:i:corollary:4-6-5-1}{4.6.5.1})
\[
  \underline{\operatorname{Lie}}(G/S)
  =
  \mathbb W\bigl(\operatorname{Lie}(G/S)\bigr);
\]
that is, for every \(S\)-scheme \(S'\),
\[
  \underline{\operatorname{Lie}}(G/S)(S')
  =
  \Gamma\!\left(
    S',
    \operatorname{Lie}(G/S)\otimes_{\mathcal O_S}\mathcal O_{S'}
  \right).
\]
By II, \SGARefLink{sga3:ii:corollary:4-1-1-1}{4.1.1.1}, the adjoint action of \(G\) gives
\(\underline{\operatorname{Lie}}(G/S)
=\mathbb W(\operatorname{Lie}(G/S))\) the structure of a
\(G\)-\(\mathcal O_S\)-module; thus
\(\operatorname{Lie}(G/S)\) is a \(G\)-\(\mathcal O_S\)-module
(cf. I, \SGARefLink{sga3:i:definition:4-7-1}{4.7.1}).  If \(G\) is moreover affine over \(S\), this says,
by I, \SGARefLink{sga3:i:proposition:4-7-2}{4.7.2}, that \(\operatorname{Lie}(G/S)\) is an
\(A(G)\)-comodule.

If \(T\) is a torus over \(S\) (IX, \SGARefLink{sga3:ix:definition:1:3}{Definition 1.3}), we say that \(T\)
is \emph{split} (``trivial'' in the original) if it is isomorphic to
\(\mathbb G_{m,S}^{\,r}\) for some integer \(r\geq 0\), and that \(T\)
is \emph{trivialized} if such an isomorphism has been fixed, or more
generally if an isomorphism
\[
  T\simeq D_S(M)
\]
has been fixed, where \(M\) is a free abelian group of rank \(r\).

Finally, recall (XII, \SGARefLink{sga3:xii:1.3}{Definition 1.3}) that a torus \(T\) of \(G\) is
called \emph{maximal} if, for every \(s\in S\), \(T_{\bar s}\) is a
maximal torus of \(G_{\bar s}\).
\SGARefTarget{sga3:xix:notations:2:5:0}

\par\medskip\noindent
\SGAInlineRefTarget{sga3:xix:theorem:2-5}\textbf{\textit{\SGARefLink{sga3:xix:theorem:2:5}{Theorem 2.5}.}}---
\textit{Let \(S\) be a scheme, let \(G\) be an \(S\)-group scheme that
is affine and of finite presentation over \(S\), with connected
fibers, and let \(s_0\) be a point of \(S\).  Suppose that \(G\) is
flat over \(S\) at \(\epsilon(s_0)\), and that the geometric fiber
\(G_{\bar s_0}\) is (reduced and) reductive
\textup{(resp. semisimple)}.  Then there exist an open subset
\(U\subset S\), containing \(s_0\), and an étale surjective morphism
\(S'\to U\) such that:}
\SGARefTarget{sga3:xix:theorem:2:5}

\begin{enumerate}
\item[\textit{\SGAInlineRefTarget{sga3:xix:theorem:2:5:item:i}(i)}]
\textit{\(G_U\) is smooth, with reductive
\textup{(resp. semisimple)} geometric fibers, and its reductive rank
and semisimple rank are constant.}

\item[\textit{\SGAInlineRefTarget{sga3:xix:theorem:2:5:item:ii}(ii)}]
\textit{\(G_{S'}\) has a split maximal torus \(T\), and the Weyl group
\textup{(cf. \SGARefLink{sga3:xix:proposition:6-3}{6.3})}}
\[
  W_{G_{S'}}(T)
  =
  \underline{\operatorname{Norm}}_{G_{S'}}(T)/
  \underline{\operatorname{Centr}}_{G_{S'}}(T)
  =
  \underline{\operatorname{Norm}}_{G_{S'}}(T)/T
\]
\textit{is finite over \(S'\).}
\end{enumerate}

